\documentclass[12pt,a4paper]{article}
\usepackage{cmap} % Makes the PDF copiable. See http://tex.stackexchange.com/a/64198/25761
\usepackage[T1]{fontenc}
\usepackage[brazil]{babel}
\usepackage[utf8]{inputenc}
\usepackage{amsmath}
\usepackage{amsfonts}
\usepackage{amssymb}
\usepackage{amsthm}
\usepackage{textcomp} % \degree
\usepackage{gensymb} % \degree
\usepackage[usenames,svgnames,dvipsnames]{xcolor}
\usepackage{hyperref}
\usepackage{multicol}
\usepackage{graphicx}
\usepackage[margin=2cm]{geometry}

\hypersetup{
    colorlinks = true,
    allcolors = {blue}
}

\newcommand*\sen{\operatorname{sen}}
\newcommand*\abs[1]{\left|#1\right|}
\newcommand*\dom[1]{\operatorname{Dom}\left(#1\right)}
\newcommand*\im[1]{\operatorname{Im}\left(#1\right)}

\newcommand*\R{\mathbb{R}}

\newcommand*\tipo{Prova I}
\newcommand*\turma{TADS121-01C}
\newcommand*\disciplina{CDI0001}
\newcommand*\eu{Helder G. G. de Lima}
\newcommand*\data{09/04/2016}

\author{\eu}
\title{\tipo - \disciplina}
\date{\data}

\begin{document}
\thispagestyle{empty}
\newgeometry{margin=2cm,bottom=0.5cm}
\begin{center}
\includegraphics[width=9.0cm]{marca}
\noindent\begin{tabular}{l c c r}
  \textbf{\disciplina (\turma)}
& \textbf{\tipo}
& \textbf{\data}
\end{tabular}
\\ Prof. \eu\footnote{
Este é um material de acesso livre distribuído sob os termos da licença \href{https://creativecommons.org/licenses/by-sa/4.0/deed.pt_BR}{Creative Commons BY-SA 4.0}.}
\end{center}

\noindent Nome do(a) aluno(a): \rule{13cm}{0.01cm}

%\section*{Instruções}
\begin{center}\fbox{
\begin{minipage}{14cm}

{\footnotesize
\begin{itemize}
\renewcommand{\theenumi}{\Roman{enumi}}
\item Identifique-se em todas as folhas.
\item Mantenha o celular e os demais equipamentos eletrônicos desligados durante a prova.
\item Anule uma das 6 questões (apenas 5 serão corrigidas). \textsc{\textbf{Questão anulada:}} \framebox(30,12){}
\end{itemize}
}

\end{minipage}
}
\end{center}


%\section*{Questões}

\begin{enumerate}
\item (2,0) Obtenha uma fórmula explícita (e simplificada) para a função $g(x) = f( f(x) )$ nos casos a seguir, e expresse o domínio e a imagem de $g(x)$ como conjuntos ou intervalos:
\begin{multicols}{2}
\begin{enumerate}
\item $f(x) = \sqrt{ -18 + x^2}$
\item $f(x) = \begin{cases}
\frac{7}{2x},& \text{ se } x \neq 0\\
0, & \text{ caso contrário.}
\end{cases}$
\end{enumerate}
\end{multicols}

\item (2,0) Explique as afirmações que são sempre \textbf{verdadeiras}, e mostre que as demais são \textbf{falsas} com contraexemplos (a argumentação lógica e as contas fazem parte da resposta):
\begin{enumerate}
\item (1,0) A função $r: \mathbb{R} \to \mathbb{R}$ definida por $\displaystyle r(x) = \sqrt[3]{8 - x^5}$ é bijetora
\item (0,5) Se $\lim\limits_{x\to a} f(x) = 0$ então com certeza $\lim\limits_{x\to a} [ f(x) \cdot g(x) ] = 0$, para qualquer $g(x)$
\item (0,5) Uma função que é par também pode ser ímpar
\end{enumerate}

\item (2,0) \textbf{Complete (e justifique suas respostas):} Sabe-se que se $p(x) = x^2 - 6x + 2$ então $\lim_{x \to 3} p(x) = -7$. Neste caso, para que os valores da função $p(x)$ fiquem a menos de $\varepsilon = 4$ unidades de distância do limite $L = -7$, é suficiente que a distância de $x$ até o ponto $a = 3$ seja menor do que $\delta = \framebox(30,15){}$ unidades de distância. E se em vez de $\varepsilon = 4$ for considerada uma distância $\varepsilon > 0$ arbitrária, basta escolher $\delta = \framebox(26,15){}$ (em função de $\varepsilon$).

\item (2,0) Calcule os seguintes limites, se existirem (ou explique porque não existem):
\begin{enumerate}
\item (1,0) $\lim\limits_{x\to +\infty} e^{[\sen(x)-x]}$
\item (1,0) $\lim\limits_{x\to 0} \dfrac{\tan(14x) - 9x}{x}$
\end{enumerate}

\item (2,0) Calcule os seguintes limites, se existirem (ou explique porque não existem):
\begin{enumerate}
\item (1,0) $\lim\limits_{x\to \pi/4} \dfrac{\tan(x-\pi/4)}{(x-\pi/4)\cos(x)}$
\item (1,0) $\lim\limits_{x\to +\infty} \tanh(x)$, sendo $\tanh(x) = \dfrac{e^x-e^{-x}}{e^x+e^{-x}}$
\end{enumerate}

\item (2,0) Se $f(x) = \dfrac{k x^2 - k}{\abs{x-1}}$ para $x \neq 1$, e $f$ é contínua em $x=1$, quanto valem $k$ e $f(1)$?

\end{enumerate}

\begin{center}
BOA PROVA!
\end{center}

\newpage
\restoregeometry
\section*{Respostas e observações}
\begin{enumerate}
\item \textit{Obtenha uma fórmula explícita (e simplificada) para a função $g(x) = f( f(x) )$ nos casos a seguir, e expresse o domínio e a imagem de $g(x)$ como conjuntos ou intervalos:}
\begin{enumerate}
\item $f(x) = \sqrt{ -18 + x^2 }$

\begin{enumerate}
\item \begin{align*}
g(x)
& = f\left( f(x) \right)
  = f\left( \sqrt{ -18 + x^2 } \right)
  = \sqrt{ -18 + \left( \sqrt{ -18 + x^2 } \right)^2 } \\
& = \sqrt{ -18 + ( -18 + x^2 ) }
  = \sqrt{ -36 + x^2 }
\end{align*}
\item \begin{align*}
\dom{g}
& = \{ x \in \R \mid x \in \dom{f} \text{ e } f(x) \in \dom{f} \} \\
& = \{ x \in \R \mid (-18 + x^2) \geq 0 \text{ e } (-36 + x^2) \geq 0 \} \\
& = \{ x \in \R \mid x^2 \geq 18 \text{ e } x^2 \geq 36 \} \\
& = \{ x \in \R \mid x^2 \geq 36 \} \\
& = \{ x \in \R \mid x \leq -6 \text{ ou } x \geq 6\} = (-\infty, -6) \cup (6, +\infty)
\end{align*}

\item $y \in \im{g}$ se existir algum $x \in \dom{g}$ tal que $g(x) = y$. Mas
\begin{align*}
y = g(x)
& \Leftrightarrow y = \sqrt{ -36 + x^2 } \geq 0 \\
& \Leftrightarrow y^2 = -36 + x^2 \\
& \Leftrightarrow y^2 + 36 = x^2 \\
& \Leftrightarrow x = \sqrt{ y^2 + 36 },
\end{align*}
então todo $y \geq 0$ esta na imagem de $g$, pois $y = g( \sqrt{ y^2 + 36 } )$, isto é, $\im{g} = \{ x \in \R \mid x \geq 0 \} = [0, +\infty)$.
\end{enumerate}


\item $f(x) = \begin{cases}
\frac{7}{2x},& \text{ se } x \neq 0\\
0, & \text{ caso contrário.}
\end{cases}$

\begin{itemize}
\item Para $x \neq 0$,
$\displaystyle g(x) = f(f(x))
= \frac{7}{2f(x)}
= \frac{7}{2\frac{7}{2x}}
= x
$
\item Para $x = 0$,
$g(x) = f(f(0))
= f(0)
= 0
= x$
\end{itemize}

Logo, para todo $x \in \R$, $g(x) = x$ e consequentemente $\dom{g} = \im{g} = \R$.

\end{enumerate}

\item \textit{Explique as afirmações que são sempre \textbf{verdadeiras}, e mostre que as demais são \textbf{falsas} com contraexemplos (a argumentação lógica e as contas fazem parte da resposta):}


\begin{enumerate}
\item \textit{A função $r: \mathbb{R} \to \mathbb{R}$ definida por $\displaystyle r(x) = \sqrt[3]{8 - x^5}$ é bijetora}

\textbf{Verdadeiro}, pois $q$ é injetora e bijetora:
\begin{itemize}
\item Injetividade:
\begin{align*}
r(a) = r(b)
& \Rightarrow \sqrt[3]{8 - a^5} = \sqrt[3]{8 - b^5}
  \Rightarrow 8 - a^5 = 8 - b^5
  \Rightarrow - a^5 = - b^5 \\
& \Rightarrow a^5 = b^5
  \Rightarrow a = b.
\end{align*}
\item Sobrejetividade:
\[
y = r(x)
\Leftrightarrow y = \sqrt[3]{8 - x^5}
\Leftrightarrow y^3 = 8 - x^5
\Leftrightarrow x^5 = 8 - y^3
\Leftrightarrow x = \sqrt[5]{8 - y^3}
\]
Consequentemente, dado qualquer $y \in \mathbb{R}$, existe $x = \sqrt[5]{8 - y^3}$ tal que
\[
r(x) = \sqrt[3]{8 - x^5}
     = \sqrt[3]{8 - (\sqrt[5]{8 - y^3})^5}
     = \sqrt[3]{8 - (8 - y^3)}
     = \sqrt[3]{y^3}
     = y.
\]

\end{itemize}

\item \textit{Se $\lim\limits_{x\to a} f(x) = 0$ então com certeza $\lim\limits_{x\to a} [ f(x) \cdot g(x) ] = 0$, para qualquer $g(x)$}

\textbf{Falso}. Por exemplo, se $f(x) = x$ e $g(x) = 1/x^3$ então $\lim\limits_{x\to 0} f(x) = 0$ mas
\[
\lim\limits_{x\to 0} [ f(x) \cdot g(x) ]
= \lim\limits_{x \to 0} x \cdot \frac{1}{x^3}
= \lim\limits_{x \to 0} \frac{1}{x^2}
= +\infty \neq 0.
\]

\item  \textit{Uma função que é par também pode ser ímpar}

\textbf{Verdadeiro}, pois a função nula é par e ímpar: de fato, se $g(x) = 0$ para todo $x$, então $g(-x) = 0 = -0 = g(x) = -g(x)$.

Na verdade, somente uma função nula pode ser par e ímpar, pois se $f$ é par então para todo $x \in \operatorname{dom}{f}$ tem-se $f(-x) = f(x)$. Se, além disso, $f$ também é ímpar, resulta que $f(-x) = -f(x)$. Consequentemente, $f(x) = f(-x) = -f(x)$, ou seja, $f(x) = -f(x)$. Mas isto implica que $2 f(x) = 0$ e que portanto $f(x) = 0$.
\end{enumerate}


\item \textit{ Sabe-se que se $p(x) = x^2 - 6x + 2$ então $\lim_{x \to 3} p(x) = -7$. Neste caso, para que os valores da função $p(x)$ fiquem a menos de $\varepsilon = 4$ unidades de distância do limite $L = -7$, é suficiente que a distância de $x$ até o ponto $a = 3$ seja menor do que $\delta = \framebox(30,15){2}$ unidades de distância. E se em vez de $\varepsilon = 4$ for considerada uma distância $\varepsilon > 0$ arbitrária, basta escolher $\delta = \framebox(26,15){ $\sqrt{\varepsilon}$ }$ (em função de $\varepsilon$). }

Para todo $x \in \R$ tem-se
\[
\abs{p(x) - (-7)}
= \abs{x^2-6x+2 +7}
= \abs{x^2-6x+9}
= \abs{x-3}^2
\]
Então, dado $\varepsilon > 0$,
\[
\abs{p(x) - (-7)} < \varepsilon
\Leftrightarrow
\abs{x-3}^2 < \varepsilon
\Leftrightarrow
\abs{x-3} < \sqrt{ \varepsilon }
\]
Em particular, se $\varepsilon = 4$ pode-se escolher $\delta = \sqrt{4} = 2$ e será verdade que $\abs{p(x) - (-7)} < \varepsilon$ quando $0 < \abs{x-3} < \delta$.

\item \textit{Calcule, caso existam:}
\begin{enumerate}

\item Observe que $e^{[\sen(x)-x]} = \frac{ e^{\sen(x)} }{e^x}$. Além disso, como $-1 \leq \sen(x) \leq 1$, conclui-se que $e^{-1} \leq e^{\sen(x)} \leq e^1$, ou seja, $e^{\sen(x)}$ é limitada. Assim, como $\lim\limits_{x\to +\infty} \frac{ 1 }{e^x} = 0$, tem-se:
\[
\lim\limits_{x\to +\infty} e^{[\sen(x)-x]}
= \lim\limits_{x\to +\infty} \frac{ 1 }{e^x} \cdot e^{\sen(x)}
= 0.
\]

\item $\lim\limits_{x\to 0} \dfrac{\tan(14x) - 9x}{x}$

\begin{align*}
  \lim\limits_{x\to 0} \dfrac{\tan(14x) - 9x}{x}
& = \lim\limits_{x\to 0} \dfrac{\tan(14x)}{x} - \frac{9x}{x}
  = \lim\limits_{x\to 0} \dfrac{\frac{\sen(14x)}{\cos(14x)}}{x} - 9
  = \lim\limits_{x\to 0} \frac{\sen(14x)}{\cos(14x)} \dfrac{1}{x} - 9 \\
& = \lim\limits_{x\to 0} 14 \frac{\sen(14x)}{14 x} \dfrac{1}{\cos(14x)} - 9
  = \lim\limits_{u\to 0} 14 \frac{\sen(u)}{u} \dfrac{1}{\cos(u)} - 9 \\
& = 14 \cdot 1 \cdot \dfrac{1}{1} - 9
  = 5
\end{align*}

\end{enumerate}

\item
\begin{enumerate}
\item \begin{align*}
    \lim\limits_{x\to \pi/4} \dfrac{\tan(x-\pi/4)}{(x-\pi/4)\cos(x)}
& = \lim\limits_{u\to 0} \frac{\tan(u)}{u}\frac{1}{\cos(u+\pi/4)}
  = \lim\limits_{u\to 0} \frac{\sen(u)}{u \cdot \cos(u)}\frac{1}{\cos(u+\pi/4)} \\
& = \lim\limits_{u\to 0} \frac{\sen(u)}{u}\frac{1}{\cdot \cos(u)}\frac{1}{\cos(u+\pi/4)}
  = 1 \cdot \frac{1}{ \cos(0) } \frac{1}{\cos(0 + \pi/4)} \\
& = \frac{1}{\cos(\pi/4)}
  = \frac{1}{ \frac{\sqrt{2}}{2} }
  = \frac{\sqrt{2}}{2}
  = \sqrt{2}.
\end{align*}

\item \begin{align*}
\lim\limits_{x\to +\infty} \tanh(x)
& = \lim\limits_{x\to +\infty} \frac{e^x-e^{-x}}{e^x+e^{-x}}
  = \lim\limits_{x\to +\infty} \frac{u-u^{-1}}{u+u^{-1}}
  = \lim\limits_{u\to +\infty} \frac{u}{u} \cdot \frac{1-u^{-2}}{1+u^{-2}}
  = \lim\limits_{u\to +\infty} \frac{1-\frac{1}{u^2}}{1+\frac{1}{u^2}} \\
& = \frac{\lim\limits_{u\to +\infty} 1-\frac{1}{u^2}}{\lim\limits_{u\to +\infty} 1+\frac{1}{u^2}}
  = \frac{1 - 0}{1 + 0} = 1.
\end{align*}

\end{enumerate}

\item \textit{Se $f(x) = \dfrac{k x^2 - k}{\abs{x-1}}$ para $x \neq 1$, e $f$ é contínua em $x=1$, quanto valem $k$ e $f(1)$?}

Para $x \neq 1$ tem-se
\[
f(x) = \frac{k x^2 - k}{\abs{x-1}}
     = k \frac{x^2 - 1}{\abs{x-1}}
     = k(x+1) \frac{(x-1)}{\abs{x-1}}
= \begin{cases}
k (x+1) & \text{ se } x \geq 1 \\
-k (x+1) & \text{ se } x < 1.
\end{cases}
\]
Assim, $\lim\limits_{x \to 1^+} f(x) = k (1+1) = 2k$ e $\lim\limits_{x \to 1^-} f(x) = -k (1+1) = -2k$. Como $f$ é contínua em $x=1$, deve ocorrer $f(1) = 2k = -2k$. Portanto, $4k=0$ e $k=0$. Assim, conclui-se que $f(1) = 2k = 2\cdot 0 = 0$.

\end{enumerate}
\end{document}
